\documentclass[11pt,reqno]{amsart}

\usepackage[T1]{fontenc}
\usepackage[utf8]{inputenc}
\usepackage{lmodern}
\usepackage{amsmath,amssymb,amsthm,mathtools}
\usepackage{booktabs}
\usepackage{microtype}
\usepackage{geometry}
\usepackage{hyperref}
\usepackage{enumitem}
\geometry{margin=1in}
\hypersetup{
  colorlinks=true,
  linkcolor=blue,
  citecolor=blue,
  urlcolor=blue,
  pdftitle={A Degree-Three Spectral Refinement of Nishioka's Lower Bound},
  pdfauthor={Bruce J. Swihart},
  pdfsubject={Three-dimensional Blaschke--Lebesgue problem},
  pdfkeywords={constant width, Blaschke--Lebesgue problem, spherical harmonics, convex geometry}
}

\newtheorem{theorem}{Theorem}[section]
\newtheorem{lemma}[theorem]{Lemma}
\newtheorem{proposition}[theorem]{Proposition}
\newtheorem{corollary}[theorem]{Corollary}
\theoremstyle{remark}
\newtheorem{remark}[theorem]{Remark}

\newcommand{\Sph}{\mathbb S^2}
\newcommand{\R}{\mathbb R}
\newcommand{\Harm}{\mathcal H}
\newcommand{\Id}{\operatorname{Id}}
\newcommand{\tr}{\operatorname{tr}}
\newcommand{\Vol}{\operatorname{Vol}}
\newcommand{\op}{\mathrm{op}}
\newcommand{\nuc}{*}
\newcommand{\dd}{\,d\sigma}
\newcommand{\ip}[2]{\left\langle #1,#2\right\rangle}
\newcommand{\norm}[1]{\left\lVert #1\right\rVert}
\newcommand{\abs}[1]{\left\lvert #1\right\rvert}

\title[A degree-three spectral refinement]{A Degree-Three Spectral Refinement of Nishioka's Lower Bound for the Three-Dimensional Blaschke--Lebesgue Problem}
\author{Bruce J. Swihart}
\date{September 2, 2026}
\subjclass[2020]{52A20, 52A40, 33C55}
\keywords{Bodies of constant width, Blaschke--Lebesgue problem, spherical harmonics, convex geometry}

\begin{document}

\begin{abstract}
Let $K\subset\R^3$ be a convex body of constant width $d$. Nishioka proved
$\Vol(K)\ge (4\pi/33)d^3$ by combining a support-function formulation with a
spherical-harmonic estimate. We retain the degree-three harmonic component of
the centered support function as an additional variable and derive a sharp
quartic determinant inequality for its curvature tensor. This gives the
candidate refinement
\[
 \Vol(K)\ge
 \frac{\pi}{1914}\left(259-27\sqrt{\frac{15183}{17303}}\right)d^3
 =0.3836027047090677\ldots\,d^3.
\]
The central tensor identity is derived invariantly and is also checked by two
exact symbolic implementations, exhaustive pairing enumeration, and
independent Python and R numerical tests. This is an AI-assisted research draft
seeking independent mathematical verification. A recent public HYRA manuscript
claims substantially stronger bounds by a different geometric and
computer-assisted method; the contribution proposed here is an independent
spectral mechanism rather than the numerically strongest current claim.
\end{abstract}

\maketitle

\begin{center}
\fbox{\begin{minipage}{0.92\textwidth}
\small\textbf{Verification status.}
This manuscript and its accompanying code form a non-peer-reviewed research
draft. Exact computer algebra supports the displayed tensor identities, but no
independent subject-matter expert has yet reviewed the complete proof. The
repository should therefore be cited as a draft seeking verification, not as a
certified or peer-reviewed theorem.
\end{minipage}}
\end{center}

\section{Introduction}

The three-dimensional Blaschke--Lebesgue problem asks for the least possible
volume of a convex body of prescribed constant width. The conjectured minimizers
are the two Meissner tetrahedra, whose common volume is approximately
$0.419860046d^3$; the problem remains open \cite{KawohlWeber2011,Nishioka2027}.

Improving Chakerian's classical estimate \cite{Chakerian1966}, Nishioka recently established
\[
 \Vol(K)\ge \frac{4\pi}{33}d^3
 =0.380799109526\ldots\,d^3
\]
using the centered support function, spherical harmonics, and Bochner's formula
\cite{Nishioka2027}. His discussion identifies a source of slack: equality in
the degree-three spectral estimate is incompatible with equality in the
pointwise curvature-matrix bound. He suggests going beyond the rotation-averaged
relaxation by introducing additional variables or symmetry-breaking information.

The present note retains the degree-three component $h_3$ of the centered
support function and estimates the curvature tensor $A_{h_3}$ more sharply than
one can estimate an arbitrary tangent-matrix field. The key new input is the
sharp inequality
\begin{equation}\label{eq:key-det-intro}
 \norm{4\det A_Y}_{L^2(\Sph)}^2
 \le \frac{15183}{17303\pi}\norm{A_Y}_{L^2(\Sph)}^4,
 \qquad Y\in\Harm_3.
\end{equation}
It quantifies a restriction intrinsic to cubic spherical harmonics and leads to
a reduced degree-three curvature budget.

Harrell's variational formulation in terms of the sum of principal radii of
curvature supplies useful historical and conceptual context: in dimension three
that formulation is a relaxation because not every admissible scalar curvature
field comes from an embedded constant-width body \cite{Harrell2002,Nishioka2027}.
The argument here instead remains within Nishioka's exact support-function
constraints.

A public HYRA manuscript dated August 2026 claims the stronger analytic bound
$0.4041207788d^3$ and a computer-assisted bound above $0.4110404737d^3$
\cite{HYRA2026}. Those claims have not been checked in this project. If correct,
they supersede the numerical coefficient obtained here. The purpose of this
note is therefore to document a complementary spectral route in a transparent,
reproducible form.

\section{Support functions and the refined spectral split}

We follow Nishioka's normalization \cite{Nishioka2027}. Let $p$ be the support
function of a smooth body of constant width $d$ and put
\[
 h=p-\frac d2.
\]
Then $h(-u)=-h(u)$. After a translation we may also impose
$h\perp\Harm_1$. Define the self-adjoint tangent endomorphism
\[
 A_h:=\nabla^2h+h\Id.
\]
The principal radii of curvature of $p$ are the eigenvalues of
$A_h+(d/2)\Id$, and constant width gives the pointwise bounds
\begin{equation}\label{eq:box}
 -\frac d2\Id\preceq A_h(u)\preceq \frac d2\Id,
 \qquad u\in\Sph.
\end{equation}
Moreover,
\begin{equation}\label{eq:volume}
 \Vol(K)=\frac{\pi d^3}{6}-\frac d2 E(h),
 \qquad
 E(h):=\int_{\Sph}\left(\frac12\abs{\nabla h}^2-h^2\right)\dd.
\end{equation}

For tangent endomorphism fields $A,B$, write
\[
 \langle A,B\rangle_2:=\int_{\Sph}\tr(AB)\dd,
 \qquad \norm{A}_2^2:=\langle A,A\rangle_2.
\]
If $Y_\ell\in\Harm_\ell$ and
$\lambda_\ell=\ell(\ell+1)$, Nishioka's calculation gives
\begin{equation}\label{eq:nishioka-mode}
 E(Y_\ell)=\frac{\lambda_\ell-2}{2}\norm{Y_\ell}_{L^2}^2,
 \qquad
 \norm{A_{Y_\ell}}_2^2
 = (\lambda_\ell-1)(\lambda_\ell-2)\norm{Y_\ell}_{L^2}^2.
\end{equation}

We first record explicitly the tensor-valued orthogonality needed below.

\begin{lemma}[Polarized Bochner identity]\label{lem:polarized-bochner}
For smooth real-valued $f,g$ on $\Sph$,
\[
 \int_{\Sph}\ip{\nabla^2f}{\nabla^2g}\dd
 =\int_{\Sph}(\Delta f)(\Delta g)\dd
 -\int_{\Sph}\ip{\nabla f}{\nabla g}\dd.
\]
\end{lemma}

\begin{proof}
Polarize the integrated Bochner identity on the unit sphere,
\[
 \int_{\Sph}\abs{\nabla^2q}^2\dd
 =\int_{\Sph}(\Delta q)^2\dd
 -\int_{\Sph}\abs{\nabla q}^2\dd,
\]
by applying it to $q=f+g$ and $q=f-g$.
\end{proof}

\begin{lemma}[Curvature-tensor orthogonality]\label{lem:tensor-orth}
If $f\in\Harm_\ell$ and $g\in\Harm_m$ with $\ell\ne m$, then
$\langle A_f,A_g\rangle_2=0$. If $f,g\in\Harm_\ell$, then
\[
 \langle A_f,A_g\rangle_2
 = (\lambda_\ell-1)(\lambda_\ell-2)
   \int_{\Sph}fg\dd.
\]
\end{lemma}

\begin{proof}
Expanding and then using Lemma~\ref{lem:polarized-bochner},
\begin{align*}
 \langle A_f,A_g\rangle_2
 ={}&\int_{\Sph}(\Delta f)(\Delta g)\dd
 -\int_{\Sph}\ip{\nabla f}{\nabla g}\dd\\
 &+\int_{\Sph}g\Delta f\dd
 +\int_{\Sph}f\Delta g\dd
 +2\int_{\Sph}fg\dd.
\end{align*}
Distinct harmonic degrees are $L^2$-orthogonal, and every term above is a
scalar multiple of $\int fg$. For equal degrees, substitute
$\Delta f=-\lambda_\ell f$ and $\Delta g=-\lambda_\ell g$ to obtain the stated
factor.
\end{proof}

Write the odd harmonic expansion as
\[
 h=h_3+h_{\ge5},
 \qquad
 h_{\ge5}=\sum_{\substack{\ell\ge5\\ \ell\ \mathrm{odd}}}h_\ell,
\]
and put
\[
 C:=A_{h_3},\qquad N:=\norm{A_h}_2^2,\qquad N_3:=\norm{C}_2^2.
\]
Since $h$ is smooth, its harmonic expansion converges in $H^2(\Sph)$; hence
Lemma~\ref{lem:tensor-orth} gives
\begin{equation}\label{eq:projection}
 N=\sum_{\substack{\ell\ge3\\ \ell\ \mathrm{odd}}}
   \norm{A_{h_\ell}}_2^2,
 \qquad
 N_3=\langle A_h,C\rangle_2.
\end{equation}

\begin{proposition}[Refined spectral split]\label{prop:spectral-split}
For every smooth admissible $h$,
\begin{equation}\label{eq:spectral-split}
 E(h)\le \frac1{58}N+\frac9{319}N_3.
\end{equation}
\end{proposition}

\begin{proof}
Equation~\eqref{eq:nishioka-mode} yields
\[
 E(h_\ell)=\frac{1}{2(\lambda_\ell-1)}
 \norm{A_{h_\ell}}_2^2.
\]
The coefficient equals $1/22$ at $\ell=3$ and is at most $1/58$ for every odd
$\ell\ge5$. Orthogonality therefore gives
\[
 E(h)\le \frac1{22}N_3+\frac1{58}(N-N_3)
 =\frac1{58}N+\frac9{319}N_3.
\]
\end{proof}

\section{The degree-three determinant inequality}

Every $Y\in\Harm_3$ is the restriction to $\Sph$ of a homogeneous harmonic
cubic
\[
 F(x)=T_{ijk}x_ix_jx_k,
\]
where $T$ is a real symmetric trace-free rank-three tensor:
$T_{ijk}=T_{(ijk)}$ and $T_{iik}=0$. Fix $u\in\Sph$ and define
\[
 P=I-u\otimes u,
 \qquad
 M(u)_{jk}=T_{ijk}u_i,
 \qquad
 v(u)=M(u)u,
 \qquad
 Y(u)=u^TM(u)u.
\]
Also put
\begin{equation}\label{eq:invariants}
 \tau=T_{ijk}T_{ijk},
 \qquad
 B_{ij}=T_{ikl}T_{jkl},
 \qquad
 B_0=B-\frac{\tau}{3}I.
\end{equation}

\begin{lemma}[Spherical curvature tensor]\label{lem:curvature-formula}
For $C=A_Y=\nabla^2_{\Sph}Y+Y\Id$, regarded as a $3\times3$ symmetric matrix
that vanishes on the normal direction $u$,
\begin{equation}\label{eq:C-formula}
 C=6PMP-2YP.
\end{equation}
Moreover,
\begin{align}
 \tr C&=-10Y,\label{eq:trace-C}\\
 \abs{C}^2&=36\abs{M}^2-72\abs{v}^2+68Y^2,\label{eq:norm-C-point}\\
 4\det C&=64Y^2+144\abs{v}^2-72\abs{M}^2.\label{eq:q-compact}
\end{align}
The determinant is taken on $T_u\Sph$.
\end{lemma}

\begin{proof}
For tangent vectors $\xi,\eta\in T_u\Sph$, the restriction formula for the
Hessian of a homogeneous cubic gives
\[
 \nabla^2_{\Sph}Y(\xi,\eta)
 =D^2F(u)[\xi,\eta]-3Y(u)\ip{\xi}{\eta}.
\]
Since $D^2F(u)=6M(u)$, adding $Y\Id$ yields
\eqref{eq:C-formula}. Trace-freeness gives $\tr M=0$ and
$\tr(PMP)=-Y$, proving \eqref{eq:trace-C}.

In an orthonormal basis whose first vector is $u$, write
\[
 M=\begin{pmatrix}Y&w^T\\ w&N\end{pmatrix},
 \qquad
 v=\begin{pmatrix}Y\\w\end{pmatrix}.
\]
Then $\abs{PMP}^2=\abs{M}^2-2\abs{v}^2+Y^2$. Expanding
$\abs{6PMP-2YP}^2$ gives \eqref{eq:norm-C-point}. Finally, for the two tangent
eigenvalues of $C$,
\[
 4\det C=2\bigl((\tr C)^2-\tr(C^2)\bigr),
\]
and substitution of \eqref{eq:trace-C}--\eqref{eq:norm-C-point} yields
\eqref{eq:q-compact}.
\end{proof}

We use the standard even moment formula
\begin{equation}\label{eq:sphere-moment}
 \int_{\Sph}u_{i_1}\cdots u_{i_{2m}}\dd
 =\frac{4\pi}{(2m+1)!!}
   \sum_{\mathcal P}\prod_{\{a,b\}\in\mathcal P}\delta_{i_ai_b},
\end{equation}
where the sum runs over all pairings of $\{1,\ldots,2m\}$. Odd moments vanish.
The quadratic consequences needed below are
\begin{equation}\label{eq:quadratic-moments}
 \int_{\Sph}Y^2\dd=\frac{8\pi}{35}\tau,
 \qquad
 \int_{\Sph}\abs{v}^2\dd=\frac{8\pi}{15}\tau,
 \qquad
 \int_{\Sph}\abs{M}^2\dd=\frac{4\pi}{3}\tau.
\end{equation}
Combining \eqref{eq:norm-C-point} with \eqref{eq:quadratic-moments} gives
\begin{equation}\label{eq:C-norm-invariant}
 \norm{C}_2^2=\frac{176\pi}{7}\tau.
\end{equation}

For the quartic calculation, define
\[
 J:=\tr(B^2),
 \qquad
 \mathcal K:=T_{abc}T_{ade}T_{bdf}T_{cef}.
\]
The next relation is special to dimension three.

\begin{lemma}[Quartic contraction relation]\label{lem:quartic-relation}
\begin{equation}\label{eq:quartic-relation}
 \mathcal K=\frac12\tau^2-J.
\end{equation}
\end{lemma}

\begin{proof}
For each $i$, let $(A_i)_{jk}=T_{ijk}$. Each $A_i$ is a symmetric traceless
$3\times3$ matrix. One has
\[
 \tau=\sum_i\tr(A_i^2),
 \qquad
 J=\sum_{i,j}\bigl(\tr(A_iA_j)\bigr)^2,
 \qquad
 \sum_iA_i^2=B.
\]
For every traceless $3\times3$ matrix $X$, Cayley--Hamilton implies
$\tr(X^4)=\tfrac12(\tr(X^2))^2$. Apply this to $X=A+tD$ and compare the
coefficient of $t^2$:
\[
 4\tr(A^2D^2)+2\tr(ADAD)
 =\tr(A^2)\tr(D^2)+2(\tr(AD))^2.
\]
Set $A=A_i$, $D=A_j$, and sum over $i,j$. The first term on the left sums to
$4J$, while
\[
 \sum_{i,j}\tr(A_iA_jA_iA_j)
 =T_{iab}T_{jbc}T_{icd}T_{jda}=\mathcal K
\]
after relabeling dummy indices and using full symmetry of $T$. The right side
sums to $\tau^2+2J$. Thus $4J+2\mathcal K=\tau^2+2J$, which is
\eqref{eq:quartic-relation}.
\end{proof}

The pairing expansion from \eqref{eq:sphere-moment}, together with
Lemma~\ref{lem:quartic-relation}, yields the six moment identities in
Table~\ref{tab:moments}. Appendix~\ref{app:pairings} gives the complete
contraction count.

\begin{table}[ht]
\centering
\caption{Quartic moments after reducing $\mathcal K$ by
\eqref{eq:quartic-relation}.}
\label{tab:moments}
\begin{tabular}{@{}lc@{}}
\toprule
Integrand & Integral divided by $\pi$ \\
\midrule
$Y^4$ & $\dfrac{16}{715}\tau^2+\dfrac{96}{5005}J$ \\[1.2ex]
$Y^2\abs{v}^2$ & $\dfrac{16}{495}\tau^2+\dfrac{32}{1155}J$ \\[1.2ex]
$Y^2\abs{M}^2$ & $\dfrac{8}{315}\tau^2+\dfrac{16}{105}J$ \\[1.2ex]
$\abs{v}^4$ & $\dfrac{16}{315}\tau^2+\dfrac{32}{315}J$ \\[1.2ex]
$\abs{v}^2\abs{M}^2$ & $\dfrac{8}{105}\tau^2+\dfrac{32}{105}J$ \\[1.2ex]
$\abs{M}^4$ & $\dfrac{4}{15}\tau^2+\dfrac{8}{15}J$ \\
\bottomrule
\end{tabular}
\end{table}

\begin{proposition}[Exact cubic determinant identity]\label{prop:det-identity}
For every $Y\in\Harm_3$ and $C=A_Y$,
\begin{equation}\label{eq:det-identity}
 \norm{4\det C}_{L^2(\Sph)}^2
 =\frac{\pi}{1001}
  \left(555264\tau^2-2265600\abs{B_0}^2\right).
\end{equation}
Consequently,
\begin{equation}\label{eq:det-bound}
 \norm{4\det C}_{L^2(\Sph)}^2
 \le \frac{15183}{17303\pi}\norm{C}_2^4.
\end{equation}
Equality in \eqref{eq:det-bound} holds exactly when $B_0=0$; this occurs, for
example, for $Y(x,y,z)=xyz$ up to scaling and rotation.
\end{proposition}

\begin{proof}
Put $a=Y^2$, $b=\abs{v}^2$, and $c=\abs{M}^2$. By
\eqref{eq:q-compact},
\[
 (4\det C)^2
 =4096a^2+18432ab-9216ac+20736b^2-20736bc+5184c^2.
\]
Substitution of Table~\ref{tab:moments} gives
\[
 \int_{\Sph}(4\det C)^2\dd
 =\frac{\pi}{1001}\left(1310464\tau^2-2265600J\right).
\]
Since $\tr B=\tau$ and $B_0=B-(\tau/3)I$,
$J=\abs{B_0}^2+\tau^2/3$, which proves \eqref{eq:det-identity}.
Dropping the nonpositive $B_0$ term and eliminating $\tau$ using
\eqref{eq:C-norm-invariant} yields
\[
 \frac{555264\pi/1001}{(176\pi/7)^2}
 =\frac{15183}{17303\pi},
\]
proving \eqref{eq:det-bound}. Equality is equivalent to $B_0=0$. For
$Y=xyz$, the six components $T_{123}$ and their permutations are $1/6$,
which gives $\tau=1/6$ and $B=(1/18)I=(\tau/3)I$.
\end{proof}

\section{The degree-three curvature budget}

Set
\[
 s:=\sqrt{\frac{15183}{17303}}.
\]
For a self-adjoint endomorphism $M$ of a two-dimensional Euclidean space,
write $\norm{M}_{\op}$ for the operator norm and $\norm{M}_{\nuc}$ for the
nuclear norm.

\begin{lemma}[Integrated nuclear-norm bound]\label{lem:nuclear}
For $C=A_Y$ with $Y\in\Harm_3$,
\begin{align}
 \int_{\Sph}\norm{C}_{\nuc}^2\dd&\le(1+s)\norm{C}_2^2,
 \label{eq:nuclear2}\\
 \left(\int_{\Sph}\norm{C}_{\nuc}\dd\right)^2
 &\le4\pi(1+s)\norm{C}_2^2.
 \label{eq:nuclear1}
\end{align}
\end{lemma}

\begin{proof}
If $\mu_1,\mu_2$ are the tangent eigenvalues of $C$, then
\[
 \norm{C}_{\nuc}^2
 = (\abs{\mu_1}+\abs{\mu_2})^2
 =\abs{C}^2+2\abs{\det C}.
\]
Writing $q=4\det C$ and integrating,
\[
 \int_{\Sph}\norm{C}_{\nuc}^2\dd
 =\norm{C}_2^2+\frac12\norm{q}_{L^1}.
\]
Cauchy--Schwarz and \eqref{eq:det-bound} give
\[
 \norm{q}_{L^1}
 \le 2\sqrt\pi\norm{q}_{L^2}
 \le2s\norm{C}_2^2,
\]
which proves \eqref{eq:nuclear2}. A second Cauchy--Schwarz inequality on the
sphere of area $4\pi$ gives \eqref{eq:nuclear1}.
\end{proof}

\begin{proposition}[Cubic curvature budget]\label{prop:cubic-budget}
For the degree-three component $C=A_{h_3}$ of any smooth admissible support
function,
\begin{equation}\label{eq:cubic-budget}
 N_3=\norm{C}_2^2\le\pi d^2(1+s).
\end{equation}
\end{proposition}

\begin{proof}
For self-adjoint tangent endomorphisms $A,M$,
\[
 \abs{\tr(AM)}\le\norm{A}_{\op}\norm{M}_{\nuc}.
\]
By \eqref{eq:box}, $\norm{A_h(u)}_{\op}\le d/2$. The projection identity
\eqref{eq:projection} therefore gives
\[
 N_3=\int_{\Sph}\tr(A_hC)\dd
 \le\frac d2\int_{\Sph}\norm{C}_{\nuc}\dd.
\]
If $N_3=0$ there is nothing to prove. Otherwise square, apply
\eqref{eq:nuclear1}, and divide by $N_3$:
\[
 N_3\le4\pi\left(\frac d2\right)^2(1+s)
 =\pi d^2(1+s).
\]
\end{proof}

\section{The volume bound}

\begin{theorem}\label{thm:main}
Every convex body $K\subset\R^3$ of constant width $d$ satisfies
\begin{equation}\label{eq:main-bound}
 \Vol(K)\ge
 \frac{\pi}{1914}\left(259-27\sqrt{\frac{15183}{17303}}\right)d^3
 =0.3836027047090677\ldots\,d^3.
\end{equation}
\end{theorem}

\begin{proof}
First suppose the support function is smooth. By \eqref{eq:box}, the two
eigenvalues of $A_h(u)$ lie in $[-d/2,d/2]$, so
\begin{equation}\label{eq:global-budget}
 N=\norm{A_h}_2^2\le 2\pi d^2.
\end{equation}
Combining Proposition~\ref{prop:spectral-split},
Proposition~\ref{prop:cubic-budget}, and \eqref{eq:global-budget},
\begin{align*}
 E(h)
 &\le \frac1{58}N+\frac9{319}N_3\\
 &\le \frac{2\pi d^2}{58}
  +\frac9{319}\pi d^2(1+s)\\
 &=\frac{\pi d^2}{319}(20+9s).
\end{align*}
Insert this in \eqref{eq:volume}:
\[
 \Vol(K)
 \ge\pi d^3\left(\frac16-\frac{20+9s}{638}\right)
 =\frac{\pi d^3}{1914}(259-27s).
\]
This proves \eqref{eq:main-bound} in the smooth class.

For an arbitrary constant-width body, use the smooth constant-width
approximation and Hausdorff continuity of volume recorded by Nishioka
\cite[Lemma 2]{Nishioka2027}, ultimately based on
\cite[Theorem 8.5.1]{MartiniMontejanoOliveros2019}. Apply the smooth estimate to
the approximating bodies and pass to the limit.
\end{proof}

\begin{corollary}\label{cor:improvement}
The coefficient in \eqref{eq:main-bound} is strictly larger than $4\pi/33$.
More precisely, with $s=\sqrt{15183/17303}$,
\[
 \frac{\pi}{1914}(259-27s)-\frac{4\pi}{33}
 =\frac{9\pi}{638}(1-s)
 =0.00280359518303217\ldots>0.
\]
\end{corollary}

\section{Reproducibility and disclosure}

The accompanying repository separates the mathematical derivation from the
computational checks. The exact Python/SymPy scripts:
\begin{itemize}[leftmargin=2em]
 \item construct a general seven-parameter symmetric trace-free cubic tensor;
 \item derive \eqref{eq:C-norm-invariant} and \eqref{eq:det-identity} by exact
       polynomial integration on $\Sph$;
 \item check \eqref{eq:q-compact} independently modulo
       $x^2+y^2+z^2-1$;
 \item enumerate all pairings behind Appendix~\ref{app:pairings}; and
 \item reproduce the exact and decimal constants in Theorem~\ref{thm:main}.
\end{itemize}
Independent base-R scripts repeat the pairing count and numerical tensor stress
tests. These computations are evidence against transcription and normalization
errors; they do not substitute for independent mathematical review.

\subsection*{Declaration of AI assistance}
During the preparation of this work, the author used OpenAI ChatGPT
(GPT-5.6 Pro, accessed August--September 2026) extensively to explore proof
strategies, carry out and cross-check symbolic calculations, draft exposition,
and prepare reproducibility code. The exact symbolic checks, independent Python
and R numerical tests, and their recorded outputs are preserved in the
accompanying repository. The named author is responsible for the decision to
release this draft, for accurately describing its provenance and verification
status, and for maintaining corrections. No AI system is listed as an author.
The complete proof has not yet received independent subject-matter review or
peer review.

\appendix
\section{Contraction counts for the quartic moments}\label{app:pairings}

Represent the four copies of $T$ by four labeled vertices, each with three
index slots. A contraction within a single tensor is a loop and vanishes because
$T$ is trace-free. Every surviving complete contraction is one of three graph
types:
\begin{enumerate}[leftmargin=2em]
 \item $D$: two disjoint vertex pairs joined by three parallel edges, giving
       $\tau^2$;
 \item $J$: two opposite doubled edges and two connecting single edges, giving
       $J=\tr(B^2)$;
 \item $\mathcal K$: the simple complete graph on four vertices, giving the
       tetrahedral contraction $\mathcal K$.
\end{enumerate}
The relation $\mathcal K=\tau^2/2-J$ then reduces every quartic moment to
$\tau^2$ and $J$.

Some scalar contractions are already fixed before applying the sphere-moment
formula. Table~\ref{tab:pairings} records the remaining pairing counts. The last
two columns result after eliminating $\mathcal K$.

\begin{table}[ht]
\centering
\caption{Exact pairing counts.}
\label{tab:pairings}
\small
\begin{tabular}{@{}lrrrrr@{}}
\toprule
Integrand & $n_D$ & $n_J$ & $n_{\mathcal K}$ & $n_{\tau^2}$ & $n_J'$ \\
\midrule
$Y^4$ & 108 & 1944 & 1296 & 756 & 648 \\
$Y^2\abs v^2$ & 12 & 216 & 144 & 84 & 72 \\
$Y^2\abs M^2$ & 6 & 36 & 0 & 6 & 36 \\
$\abs v^4$ & 4 & 40 & 16 & 12 & 24 \\
$\abs v^2\abs M^2$ & 2 & 8 & 0 & 2 & 8 \\
$\abs M^4$ & 1 & 2 & 0 & 1 & 2 \\
\bottomrule
\end{tabular}
\end{table}

For example, in $Y^4$ the disconnected type contributes
$3(3!)^2=108$ pairings, the tetrahedral type contributes $(3!)^4=1296$, and
the six labeled $J$ graphs contribute
$6\cdot 3^4(2!)^2=1944$. The other rows follow by the same slot-counting
argument with the scalar contractions in $\abs v^2$ and $\abs M^2$ fixed.
The repository script \texttt{python/verify\_pairing\_counts.py} exhaustively
enumerates the pairings and reproduces the table exactly; the base-R script
\texttt{R/verify\_pairing\_counts.R} is an independent implementation.

Multiplication by $4\pi/(2m+1)!!$, where $2m$ is the number of remaining
sphere-coordinate factors, gives Table~\ref{tab:moments}.

\begin{thebibliography}{99}

\bibitem{Chakerian1966}
G.~D. Chakerian,
\emph{Sets of constant width},
Pacific J. Math. \textbf{19} (1966), no.~1, 13--21.

\bibitem{Harrell2002}
E.~M. Harrell II,
\emph{A direct proof of a theorem of Blaschke and Lebesgue},
J. Geom. Anal. \textbf{12} (2002), no.~1, 81--88.

\bibitem{HYRA2026}
HYRA,
\emph{A certified geometric lower bound for the three-dimensional
Blaschke--Lebesgue problem},
public research manuscript in the Tencent Hunyuan Hyra-results repository,
dated August 15, 2026,
\url{https://github.com/Tencent-Hunyuan/Hyra-results/tree/main/AI4Science/3d_blaschke_lebesgue}.

\bibitem{KawohlWeber2011}
B.~Kawohl and C.~Weber,
\emph{Meissner's mysterious bodies},
Math. Intelligencer \textbf{33} (2011), no.~3, 94--101.

\bibitem{MartiniMontejanoOliveros2019}
H.~Martini, L.~Montejano, and D.~Oliveros,
\emph{Bodies of Constant Width: An Introduction to Convex Geometry with
Applications},
Birkh\"auser, 2019.

\bibitem{Nishioka2027}
A.~Nishioka,
\emph{An improved lower bound for the three-dimensional Blaschke--Lebesgue
problem from spectral and dual perspectives},
J. Math. Anal. Appl. \textbf{566} (2027), 131038,
\url{https://doi.org/10.1016/j.jmaa.2026.131038}.

\end{thebibliography}

\end{document}
